\chapter{FORTE FRECCIA}\label{forte-freccia}

Al complesso del Vettore, feci rapporto a un certo capitano Price,
nell'edificio amministrativo, una struttura in cui l'Unef era subentrata
ai ruhar che avevano gestito le operazioni del Vettore. Entrando dal
campo d'aviazione, vidi un sacco di ruhar fuori dalla linea di
recinzione della base Unef chiamata Forte Freccia. Dal momento che gli
esseri umani non sapevano come far funzionare o mantenere l'enorme
cannone a rotaia del Vettore che sparava il carico in orbita, il
cambiamento più grande aveva riguardato la bandiera che sventolava sopra
la base, visto che a gestire il reattore e il resto del complesso del
Vettore erano i criceti. L'Unef aveva ricavato Forte Freccia dalla città
operaia che era cresciuta vicino al reattore a fusione, prendendo il
controllo degli edifici esistenti e impiegando il minimo sforzo per
convertirli a uso militare. Non faceva una piega, l'Unef non aveva
intenzione di rimanere su Paradiso abbastanza a lungo da investire in
infrastrutture. C'era un corridoio semisicuro tra l'aeroporto e Forte
Freccia, su entrambi i lati del corridoio c'erano criceti che andavano
avanti con le loro vite.

L'assistente del capitano Price mi lasciò ad aspettare quasi un'ora,
durante la quale rimasi seduto a guardare la porta dell'ufficio. Price
era lì dentro, lo sentivo parlare al telefono e fare lunghe pause. Non
era così impegnato da non riuscire a darmi il benvenuto a bordo. Sul
Dodo non avevano servito la colazione, avevo fame e volevo scroccare del
cibo prima che la mensa chiudesse. Alla fine, il capitano Price si
presentò sulla porta del suo ufficio, mi guardò con aperta ostilità e
grugnì quando mi alzai in piedi e gli feci il saluto. «Bishop.» Era
un'asserzione.

«Sì, signore, sergente Bishop a rapporto.» Non avevo un ordine del
giorno da consegnarli, tutto era sui nostri telefoni e tablet. Nel suo
ufficio, non m'invitò a sedermi, rimasi quasi sull'attenti, a destra
della porta.

«Bishop», ripeté, indicando qualcosa che non potevo vedere sul suo
tablet. «Il tizio di Barney.»

Merda, quella faccenda mi stava davvero nauseando.

«Ho i tuoi ordini qui», continuò, «e il tuo fascicolo personale. La
mancanza di disciplina sembra essere un'abitudine per te. Come agire in
modo avventato. Non abbiamo intenzione di tollerare nulla di tutto ciò,
qui a Forte Freccia. Potresti pensare che i tuoi quindici minuti di
gloria ti diano diritto a un trattamento speciale.» Non mi preoccupai di
protestare, perché si era già fatto un'idea di me. «È un November Golf,
capito?»

«Sì, signore.» Basso profilo, mi dissi. Il maggiore Perkins, e il
quartier generale dell'Unef attraverso di lei, voleva che garantissi un
basso profilo. Tenere la bocca chiusa era il primo passo.

Un semplice ``sì'' da parte mia non andava abbastanza bene, a quanto
pare, Price aveva un diavolo per capello e un pulpito da cui predicare,
e stava per sfruttare al massimo il fatto di avere un pubblico
prigioniero. ``Avrai sentito dire che Forte Freccia è una discarica per
casinisti e malcontenti.» In realtà, non l'avevo sentito e non avevo
sentito nessuno usare la parola ``malcontento'' da quando ero, tipo,
alle elementari. «Il quartier generale dell'Unef pensa che, poiché i
ruhar hanno bisogno del Vettore per spedire i loro cereali fuori dal
pianeta, Forte Freccia sia al sicuro dagli attacchi e non abbiamo
bisogno di una presenza significativa della sicurezza qui.» Ok, non si
trattava di me, allora, Price aveva un problema con il quartier generale
dell'Unef e io ero un bersaglio conveniente per la sua contrarietà. «Si
sbagliano, l'importanza del Vettore ne fa il primo posto che i ruhar
cercheranno di occupare, se mai proveranno a riconquistare il pianeta.
Questa è una forza d'élite.» Tamburellò due volte col dito sulla
scrivania per enfatizzare, producendo un suono stridente. L'effetto che
il gesto ebbe su di me fu farmi pensare che il tizio avesse bisogno di
tagliarsi le unghie. «Élite. Noi di Forte Freccia dobbiamo mostrare ai
ruhar di essere così forti qui, che non vale la pena cercare di prendere
questo sito.»

Mi limitai ad annuire, in silenzio, perché non volevo mostrarmi in
disaccordo con la sua logica fallace. Se i ruhar fossero riusciti a
mandare in orbita una flotta abbastanza potente da scacciare i kristang,
gli umani intrappolati sulla superficie avrebbero rappresentato un
ostacolo. I criceti potevano stazionare in orbita e usare colpi di
precisione con cannoni a rotaia, maser e missili intelligenti per
eliminare la resistenza umana, a prescindere da quanto forte fosse la
guarnigione di presidio a Forte Freccia.

«Non abbiamo un posto per te qui e potrebbe servirti un po' di aiuto per
tenerti fuori dai guai. Ti assegnerò il compito di scorta ai convogli,
ti presenterai al sergente scelto Lombard domattina. Cerca di rigare
dritto, e\ldots»

``Signore?'' L'assistente di Price lo chiamò da fuori dall'ufficio. «Il
colonnello Young sta venendo qui proprio ora.»

Il colonnello Young entrò a grandi passi, mollò lo zaino su un tavolo,
si passò una mano sulla faccia per asciugarsi il sudore e infilò la
testa nella porta dell'ufficio di Price. «Porca miseria, siamo nella
merda adesso. I nostri amici kristang tireranno fuori il loro
cacciatorpediniere domani, il che fa sì che ci resti un'unica fregata in
orbita per il supporto al fuoco. Cosa ci scommetti che la fregata se la
svignerà al primo segno di una nave da guerra dei criceti? Il vecchio
vuole piani di emergenza, nel caso in cui i ruhar ritornino. Dobbiamo
essere pronti a difendere questo pianeta da soli.»

«Va così male, signore?», chiese il capitano Price, spaventato.

«Abbastanza male, perché vogliono che un plotone se ne vada da qui
domani e sia dispiegato altrove per aumentare la sicurezza in un paio di
basi logistiche. L'Unef pensa che i ruhar non rischierebbero di
danneggiare il Vettore di carico, so che ci stanno sparpagliando per
coprire obiettivi più probabili. Che?» Il colonnello Young reagì quando
Price alzò il sopracciglio. Si girò a destra e mi notò: «Chi sei,
sergente?».

Salutai in modo risoluto. «Bishop, signore, mi hanno trasferito qui,
sono arrivato questa mattina.»

«Bishop, eh? Sì, sei il tizio di Barney, avevamo sentito che stavi
arrivando.» Non ne sarei mai uscito. «Te ne puoi andare, sergente. E
tieni la bocca chiusa. Non che possiamo mantenere il segreto comunque»,
disse Young con un sospiro.

«Sì, signore.» Uscii in fretta, quanto più in fretta era possibile senza
perdere del tutto la mia dignità.

\hfill\break

La squadra di scorta ai convogli cui ero stato assegnato era un buon
gruppo, guidato da un sottotenente novellino abbastanza sveglio da dare
retta al suo sergente scelto. Il capitano Price aveva ragione, non c'era
posto per un altro sergente nella squadra di scorta, così tappai il buco
in sostanza come specialista esperto, feci quello che il sergente
Lombard mi diceva di fare, tenni la bocca chiusa e mi comportai bene,
almeno all'esterno. Indossavo ancora i gradi di sergente, portavo
un'arma da fianco e nel database del personale ero ancora nella lista
degli E5, quindi non ero stato abbassato di grado. Non ancora.

Nel complesso l'incarico non era male, ebbi modo di vedere una fetta più
grande di Paradiso, dalle giungle equatoriali, oltre passi di montagna e
giù per le praterie. Era un bel pianeta, peccato che fosse infestato dai
criceti e che avremmo dovuto consegnarlo ai kristang dopo il giorno
stabilito per l'evacuazione. Paradiso sarebbe stato una bella seconda
casa per l'umanità. Dal momento che non eravamo neppure in grado di
raggiungere i pianeti del nostro sistema solare senza l'aiuto degli
alieni, era solo un pio desiderio.

Un tipico viaggio del convoglio durava quattro o cinque giorni
all'andata, quattro o cinque giorni al ritorno, trasportando cereali e
altro cibo per criceti. I criceti civili lasciavano il pianeta con
l'ascensore spaziale, non con il Vettore, perciò non avevamo a che fare
con il caos delle famiglie di criceti infelici. I ruhar non ci causavano
problemi, a parte atti minori, fastidiosi e casuali di sabotaggio qua e
là; li interpretavo come un modo dei criceti di alzare il dito medio
all'Unef. Cosa che potevo capire, noi umani eravamo la forza occupante,
agivamo come sicari dei kristang agli occhi dei ruhar e, se fossi stato
nella loro situazione, avrei fatto il ribelle. Alcune delle famiglie
ruhar erano lì da tre generazioni, avevano messo radici nel terreno
fertile di Paradiso e avevano coltivato e badato ai propri affari in
modo pacifico, finché il recente spostamento del wormhole aveva fatto
decidere ai kristang che era giunto il momento di riprendersi il
pianeta. L'Unef richiedeva scorte ai convogli per prevenire problemi, la
preoccupazione era che i criceti avrebbero potuto ribellarsi perché ci
avvicinavamo al giorno dell'evacuazione ed era chiaro che stavano
davvero lasciando questo pianeta, con tutta probabilità per sempre. I
cinesi e i francesi avevano trovato criceti vagabondi in luoghi che si
supponeva fossero stati evacuati del tutto e truppe di ogni nazionalità
avevano scoperto depositi segreti di armi per criceti. Forse i ruhar non
avevano intenzione di andarsene in silenzio, dopotutto. Tra un viaggio
del convoglio e l'altro, trascorrevamo un giorno o due a Forte Freccia,
il che significava dormire in un vero letto, mangiare cibo caldo, che
qualcun altro cucinava per te, e poter usare la palestra, i campi da
baseball e altre opportunità di R\&r, riposo e recupero. Forte Freccia
aveva anche una grande piscina, il cui punto di forza era la possibilità
di vedere soldatesse in costume da bagno. Purtroppo, non una sola cazzo
di volta che andai nella sala mensa di Forte Freccia facevano
cheeseburger. Mangiai del buon pesce e patatine, un polpettone decente,
che era più pane che carne, e pasticcio di pollo che era più verdure e
pane che pollo. Sembrava che su Paradiso scarseggiasse qualsiasi tipo di
carne, ci erano giunte voci che le navi di rifornimento dalla Terra
fossero in ritardo, o che operassero a orari irregolari, o che le
battaglie spaziali altrove nel settore stessero inducendo i thuranin a
deviare le loro navi. Poco dopo che noi umani eravamo subentrati ai
ruhar al Vettore, il comandante della base dell'Unef aveva ordinato di
piantare un orto, così avremmo avuto verdure fresche da mangiare.
Pomodori, meloni, cipolle, spinaci, peperoni, tutto quello che si vede
in un tipico mercato contadino. Seppi che la situazione stava diventando
preoccupante un giorno che le uniche opzioni per il pranzo nella sala
mensa di Forte Freccia erano insalata di spinaci e fajita vegetariana.
Certo, gli spinaci sono una buona fonte economica di proteine e, quando
arricchii la mia insalata con noci e crostini, il tutto era abbastanza
gustoso, ma a volte un soldato vuole un pezzo di carne da masticare. E
formaggio. E un panino. Tostato. Con ketchup. Avere anche le cipolle
fritte sarebbe bello.

Mentre ci avvicinavamo a Forte Freccia alla fine del mio quinto viaggio
in convoglio, non vedevamo l'ora di trascorrere tre giorni di R\&r alla
base, perché i nostri veicoli sarebbero stati sottoposti a manutenzione
ordinaria. Il soldato Pope si chinò verso di me per parlare al di sopra
del ronzio del motore elettrico del camion. «Sergente, ha piani per il
R\&r di domani?»

Quando ero entrato a far parte dell'unità di scorta, la reazione
iniziale degli altri era stata: primo, l'inevitabile curiosità per la
mia piccola celebrità; secondo, un mucchio di domande su come fossi
riuscito a fare abbastanza casino da farmi assegnare al servizio di
scorta a Forte Freccia. Le battute su Barney le prendevo in modo
sportivo; le avevo sentite un milione di volte ormai e avevo una
risposta bonaria pronta per tutte. All'inizio, il nostro tenente era
scettico nei miei confronti, poi il sergente scelto Lombard mi aveva
lasciato gestire un po' del carico di lavoro, affidandomi faccende
amministrative minori che gli facevano perdere del tempo, perciò, quando
la gente l'aveva visto, aveva pensato che fossi un tipo a posto, ero
stato accettato ed era stato bello. «Nessun piano, perché?» Era previsto
un tempo caliginoso, caldo e umido, possibilità di rovesci pomeridiani,
tipico vicino all'equatore in quel periodo dell'anno.

«C'è un'astronave kristang che si è schiantata a tre chilometri a nord
della base», spiegò Pope. «Un gruppo dei nostri andrà a controllare, se
vuole, venga con noi.»

«Un'astronave è caduta dall'orbita? Cazzo! Ne è rimasto qualcosa?» Non
cercai di nascondere il mio entusiasmo.

«Sì, sì, ho visto foto di persone che ci sono state. Dicono che è una
fregata. Risale alla battaglia in cui i ruhar hanno preso questo posto
ai kristang l'ultima volta, quindi è vecchia, e la giungla ne ha
inghiottita una parte, ma la struttura principale è ancora lì. I criceti
hanno ripulito le armi e il reattore. L'Unef ci dissuade dal curiosare
da quelle parti, perché i kristang sono riservati su questo, ma in molti
lo fanno.»

«Sì, mi, ehm, mi piacerebbe andare.» Scacciai il pensiero di prendere
dei souvenir, è probabile che i kristang si sarebbero accigliati non
poco per questo. «Grazie.»

Pianificammo di dirigerci verso la nave abbattuta la mattina presto.
Mentre facevamo colazione, il soldato Crockett ci si avvicinò in fretta
e sussurrò: «Sbrigatevi, andiamocene da qui. C'è un generale indiano che
atterrerà qui in visita, terrà un discorso subito dopo pranzo e vogliono
che tutti gli uomini qui in sala mensa gli facciano ressa attorno».

Tutti brontolammo. Nessuno aveva voglia di stare seduto nella soffocante
sala mensa, ad ascoltare un altro discorso noioso. Trangugiammo toast e
uova in polvere, afferrammo i panini al burro di arachidi per il pranzo
e attraversammo la base in pratica di corsa, per scomparire nella
giungla il più presto possibile. Chiunque se ne fosse stato seduto senza
granché da fare stava per essere cooptato come volontario obtorto collo
per scaldare sedie in sala mensa.

Tre chilometri dall'astronave abbattuta si rivelarono essere più di
sedici, e non c'era una strada, perciò andammo a piedi. Quella via era
stata percorsa da un numero di persone sufficiente a tracciare una sorta
di sentiero attraverso la giungla. Dico una sorta, perché il sentiero si
era evoluto a mano a mano che ogni gruppo di persone aveva trovato modi
migliori per arrivare lì; aggirare le colline, trovare punti meno
profondi per attraversare i torrenti, evitare fango e densi fossi di
spine. Sarebbe stato meglio se il sentiero fosse stato segnato, ma non
lo era. Qualunque tipo di animale nativo vivesse in quella giungla aveva
creato sentieri e c'erano un sacco di vicoli ciechi. In sostituzione dei
segni, cercammo di seguire qualsiasi traccia sembrasse più trafficata,
ma questo si rivelò una cattiva idea, perché le persone passate prima di
noi erano degli idioti. Più di una volta affondai fino alle ginocchia
nei fangosi ruscelli della giungla. Era caldo e umido e c'erano grandi
insetti inquietanti, anche se sapevamo, o meglio ci avevano detto, che
il loro veleno non poteva colpirci. Continuavo a schiacciare e a
strizzare insetti che cadevano dagli alberi e mi atterravano sulla nuca,
sperando in un pasto facile. È chiaro che non avevano ricevuto il
promemoria in cui si diceva che gli umani non erano commestibili per gli
organismi biologici nativi di Paradiso. Oppure erano solo degli odiosi
figli di puttana e volevano mordere o pungere qualcosa.

Quella possibilità di rovesci pomeridiani, che ai tropici equivale a
``possibilità che il sole sorga al mattino'', si rivelò essere la
valanga battente di un nubifragio. Durò meno di cinque minuti, che
sembrarono molto più lunghi, visto che al primo minuto ci ritrovammo
fradici fino all'osso. Alcuni di noi cercarono di rannicchiarsi sotto
gli alberi più grandi a foglia larga, finché non iniziarono a cadere
fulmini e tutti ci allontanammo il più possibile da quel riparo. Quando
la pioggia cessò, uscì il sole, ed era come camminare attraverso un
bagno turco. Grandi e grosse gocce d'acqua riscaldate dal calore solare
cascavano dagli alberi sulle nostre teste, l'aria si tagliava con un
coltello e tutti gli insetti che prima della tempesta stavano dormendo
adesso erano svegli e affamati. Era, concordammo tutti, molto meglio che
ascoltare un discorso in sala mensa.

Non avrei mai trovato la nave da solo. Tutti noi umani facevamo troppo
affidamento sulla tecnologia avanzata, anche sulla Terra, e avevo
iniziato a perdere abilità di base come la lettura delle mappe e la
navigazione sul campo. Infine, ci imbattemmo nel mezzo distrutto, per lo
più seguendo una scia d'involucri di cibo pronto militare gettati via.
Da quello che avevo sentito, mi aspettavo che non fosse rimasto altro
che le ossa della struttura della nave, ma per la mia gioia era molto
più intatta di quanto credessi. Era enorme e, se le fregate kristang
erano di quelle dimensioni, non avrei voluto incontrare uno dei loro più
grandi combattenti nella guerra spaziale. Mancava un grosso pezzo della
sezione di poppa, immagino nel punto in cui i ruhar avevano rimosso il
reattore a fusione, e la prua era sepolta nel terreno paludoso. Per quel
che potevo vedere dallo spazio tra le due sezioni tra loro più lontane
della nave distrutta, stimai che era molto più grande di un sottomarino
nucleare, forse lunga quanto una portaerei. Forse ancora di più. La
maggior parte della nave era occupata dalla sala macchine, non sapevo se
quella parte fosse stata pressurizzata con aria respirabile, ma
immaginavo che per praticità fosse così, perché i motori dovevano avere
bisogno di manutenzione. Riuscimmo a entrare nella nave e vagare un po'
in giro con le torce, non c'erano molti animali pericolosi su Paradiso,
il che era una cosa grandiosa visto che l'unica arma che avevamo con noi
era la mia pistola da fianco. Anche l'interno era sporco, infangato e
rivestito di vegetazione, pieno di insetti. Dopo avere curiosato per un
po', iniziammo ad annoiarci e, ovviamente, ci venne fame. Qualcuno
suggerì di salire in cima alla nave, dove non saremmo stati seduti su un
terreno paludoso.

Era una bella giornata, lontano dalla base, fuori nella boscaglia,
intento a esplorare qualcosa di nuovo, senza nessuno che mi sparava
addosso. Avevo una borraccia piena di succo energetico, un panino al
burro di arachidi, una barretta Hooah! e un sacchetto di frutta secca.
Che altro si potrebbe chiedere? Si diceva che sul menu della mensa ci
fosse pollo quella sera e pensavo di farmi una nuotata nella piscina
della base, magari giocare a basket o a softball più tardi.

Schermandomi gli occhi con la mano, guardai in alto per controllare la
posizione del sole, o stella locale, o qualunque fosse il termine
corretto, e vidi la tipica luce scintillante di un'astronave che saltava
nello spazio normale. Un'altra nave da trasporto ruhar? Ne avevamo viste
con una certa regolarità. Mi affascinava ancora l'idea di navi che
viaggiano più veloci della luce.

Un altro scintillio. Niente d'insolito, le navi da trasporto ruhar
arrivavano spesso in formazione multipla ed erano sempre scortate da
navi kristang, di solito fregate, non che riuscissi a vedere la
differenza tra una nave e l'altra dalla superficie. Mmh. Altre luci
scintillanti. E altre ancora.

Molte altre.

«Ehm, ehi, ragazzi», Pope indicò il cielo: «Ci sono un sacco di navi
lassù».

La task force dei kristang era di ritorno?

Il mio zPhone emise un segnale strozzato e si zittì. Oh merda. Non
prometteva niente di buono. Tamburellai col dito sull'icona del canale
del comando, ma nulla si mosse. Stavo già cercando di contattare
chiunque, proprio chiunque, sul mio telefono, quando vidi una scia di
luce infuocata scendere dal cielo e un'enorme esplosione nella direzione
di Forte Freccia. Una fontana di terra esplose verso l'alto
all'orizzonte, evolvendosi all'istante in una nube a forma di fungo. Lo
avevamo visto tutti nei video didattici: il dart di un cannone a rotaia.
Un piccolo, denso proiettile di tungsteno, o, più probabile, di un
qualche materiale esotico alieno, accelerò fino a una percentuale
significativa di velocità della luce, scavando a fuoco un buco
nell'atmosfera e schiantandosi contro il pianeta. Dentro Forte Freccia.
Mentre guardavo impotente, a bocca spalancata, altre scie di
condensazione sfrecciavano verso il basso, per colpire la base. Quelle
scie curvavano in volo. Missili intelligenti iperveloci, che seguivano a
ruota i proiettili del cannone a rotaia. Ci furono altre esplosioni,
tutte provenienti da Forte Freccia.

Tutti ebbero la mia stessa reazione. Primo, porca puttana! Secondo, che
cazzo faccio ora? Senza ordini, scendemmo tutti a precipizio dalla nave
abbattuta. Mi guardai attorno per vedere se per caso qualcosa fosse
cambiato dall'ultima volta che avevo controllato un minuto prima, ma la
situazione era la stessa: diciassette persone e, tra tutti, avevamo
proprio un'arma soltanto, a meno di non contare i coltelli. La mia
pistola non sarebbe stata molto utile contro i ruhar.

«Sergente, cosa succede?»

Sergente. Ogni faccia era rivolta verso di me. Merda. Tutti gli altri
erano soldati o specialisti. Cazzo, ero un sergente, no? Quella mattina,
avrei dovuto essere solo uno dei tanti in escursione nei boschi. Ora i
galloni sulla giacca della mia uniforme e la mia arma da fianco
significavano che dovevo fare qualcosa. Qualsiasi cosa.

«Ne so quanto voi. Avete tutti ricevuto un segnale sullo z-Phone?» Le
persone scossero la testa in segno di diniego. Almeno tutti avevano
avuto la prontezza di controllare il telefono dopo che la base era stata
colpita. «Devono essere i ruhar che disturbano le nostre comunicazioni,
non ricevo neppure il segnale di navigazione.» Non potevo neanche
utilizzare la funzione di prossimità per vedere sulla mappa i telefoni
delle persone attorno a me. L'intero sistema doveva essere stato
disattivato. «Bene, gente, il network è fuori uso, impostate i telefoni
in modalità di sola ricezione.» Si supponeva che questo avrebbe impedito
agli zPhone di trasmettere qualsiasi segnale, anche la nostra posizione.
L'Unef sospettava che i kristang avessero un modo di localizzarci anche
se il telefono era in modalità stealth, ma non avevamo molta scelta.
L'ultima cosa di cui avevamo bisogno era che i criceti ci
localizzassero: si sperava che non potessero attingere alla tecnologia
kristang. Essendo probabile che la base fosse stata abbattuta, un
manipolo di diciassette esseri umani avrebbe rappresentato un
bell'obiettivo secondario. «Prendete la vostra attrezzatura», dissi in
automatico, ignorando che l'attrezzatura a quel punto erano zaini e
borracce. «Torniamo alla base, il prima possibile.»

«Base?» Il soldato Collins indicò la colonna di fumo e, mentre terminava
la sua domanda, ci fu un'altra esplosione in quella direzione. «Non c'è
nessuna base! Non può esserne rimasto nulla.»

«Non lo sappiamo. Ma è sicuro come la morte che non ce ne resteremo
nascosti qui nella giungla. Torneremo alla base, perché potrebbero
esserci persone che hanno bisogno del nostro aiuto e perché è il nostro
lavoro, il nostro dovere. Se avete bisogno di più motivazione, i
kristang sono il nostro unico passaggio per tornare a casa e l'unico
modo per avere approvvigionamenti di cibo via nave. Se i ruhar si
stabiliscono qui prima che tornino i kristang, siamo tutti nella merda
fino al collo.»

«Se le cazzo di lucertole tornano», brontolò Collins, «come impediremo
ai criceti di accamparsi di nuovo qui? Non abbiamo nemmeno armi.»

«Li terremo occupati e sbilanciati e li colpiremo ogni volta che
possiamo, per guadagnare tempo affinché la task force kristang torni
qui.» Guardai attorno a me quel gruppo di persone che conoscevo appena,
persone provenienti da diverse squadre di scorta al convoglio. Non
sapevo nemmeno tutti i loro nomi, ci eravamo incontrati solo alla mensa
quella mattina. «Non conosco tutti voi, con alcuni sono stato in
servizio di scorta e altri mi hanno detto di non essere venuti fin qui
per giocare a fare da bambinaie a un gruppo di criceti.» Era una
lamentela comune all'interno dell'Unef. «Questa è la nostra possibilità
di riscatto.»

«Sergente, sono tutto per il dovere», disse Pope, «ma Collins ha
ragione: cosa dovremmo fare senza armi?»

La specialista Amaro prese la parola, giuro che si alzò in punta di
piedi e sollevò la mano come fosse un'alunna della scuola elementare:
«C'è un deposito di munizioni fuori dalla base, è in un magazzino dei
criceti, costruito nel fianco della collina, lungo la strada d'accesso
che costeggia il binario del Vettore. Ci ho consegnato le scorte un mese
fa, penso che sia ancora lì».

«Presidiato? Ci sono guardie?», chiesi. In quel caso, era probabile che
non avremmo avuto modo di aprire la porta del bunker.

«Due ragazzi, il giorno che ci sono stata. C'è una grossa porta pesante
all'ingresso e una specie di piccola baracca che abbiamo messo su per le
guardie, i criceti non avevano niente del genere, l'avevano lasciata
senza presidi. È forse, non lo so\ldots», guardò la mappa sul suo
zPhone, «merda, senza la funzione Gps non so dove sia. Non lontano?
Potete vedere dove hanno tagliato la collina per fare la strada di
accesso.» Puntò il dito a ovest, indicando una cicatrice orizzontale
lungo la montagna. Soltanto una sezione del taglio era visibile
attraverso gli alberi, o qualunque cosa fossero. Avevo visto di sfuggita
la strada di accesso sul volo per Forte Freccia. Le mappe dicevano che
c'era una strada di accesso che correva lungo entrambi i lati del
Vettore, con strade secondarie che portavano al tubo di lancio del
Vettore stesso, ogni chilometro e mezzo circa, o qualcosa del genere.

«Bene, qualcun altro c'è stato? No?» Tutt'intorno la gente fece di no
con la testa. Ero tentato di salire sulla nave per avere una visuale
migliore. «Possiamo andare dritti alla strada d'accesso da qui?»

Amaro sembrava colpita: «Ehm, non lo so, sergente?». Nell'esercito,
``non lo so'' è una risposta accettabile in pieno, molto meglio che
cercare di arrampicarsi sugli specchi a suon di stronzate. Di solito
``non lo so'' dovrebbe essere seguito da ``ma lo scoprirò''. «Penso che
ci sia una scarpata tra qui e la strada\ldots{} non posso, ehm,
accidenti a questa cosa.» Tamburellò con le dita sullo schermo del suo
zPhone.

«Va tutto bene, Amaro, ci siamo tutti affidati troppo alla tecnologia
sofisticata e abbiamo lasciato che le nostre abilità di navigazione di
base si indebolissero, so di averlo fatto anche io.» Ed era vero visto
che, senza il Gps, non avevo la minima idea di dove fossimo. Strinsi le
cinghie del mio zaino e guardai le facce che mi fissavano a loro volta.
A volte, l'unica cosa di cui la gente ha bisogno è sentire che qualcuno,
chiunque sia, ha una specie di piano. «Gente, prendiamo armi e
munizioni. Amaro, fai strada, stabilisci un ritmo sostenibile, potrebbe
essere una lunga corsa.»

\hfill\break

Era una lunga corsa nel caldo, almeno cinque, forse sei chilometri, e
sembrava che un chilometro e mezzo fosse in salita. Diventò più facile
quando raggiungemmo la strada di accesso, Amaro riconobbe un declivio
roccioso e capì da che parte andare lungo la strada. Accelerammo il
passo nonostante il calore crescesse e l'acqua nelle borracce fosse
finita, sulla strada era più facile correre e avevamo l'incentivo
aggiunto di vedere le navi d'attacco ruhar ronzare sopra le nostre
teste. Ogni volta che avvistavamo uno di quegli uccelli d'attacco che
l'Unef chiamava Avvoltoi, saltavamo fuori dalla strada per metterci al
riparo. Nasconderci sotto gli alberi ci faceva sentire un po' più al
sicuro, la mia ipotesi era che, con la loro tecnologia, i ruhar
sapessero con precisione dov'eravamo. Non valeva la pena sprecare
munizioni per un piccolo gruppo di umani primitivi disarmati.

Raggiungemmo la strada secondaria che conduceva al deposito di
munizioni; avevamo corso solo qualche centinaio di metri lungo la via
d'accesso, quando da un cespuglio sul ciglio una voce gridò: «Fermi!
Restate lì!». La voce fu più vibrante di quanto sperassi di sentire,
mentre il suo proprietario mi puntava il fucile alla testa.

«Sergente Joe Bishop, X fanteria. Chi comanda qui?»

«Lei», disse un'altra voce e un tizio sbucò fuori da dietro un albero.
«Siamo solo noi, sergente. Sono lo specialista Rogen e quello sotto i
cespugli è il soldato Wayne.»

Merda. Avevo sperato dentro di me che al deposito munizioni ci fosse
qualcuno di rango più elevato per togliermi dalle spalle il peso della
responsabilità.

Rogen tamburellò con le dita sullo zPhone alla sua cintura: «Le nostre
comunicazioni sono interrotte». Notai che Rogen stava ancora puntando il
fucile nella nostra direzione, con la canna un po' abbassata, ma il dito
ben posizionato accanto al grilletto. Lui e Wayne erano in uniforme, noi
in calzoncini e maglietta a maniche corte, e loro non conoscevano il
nuovo arrivato che diceva di essere sergente. Rogen guardò al di sopra
della mia spalla: «Ehi, tu sei Miller, giusto? Sei in servizio di
scorta».

«Sì», ammise Miller. «Sei in una squadra di baseball, interbase?»

«Seconda base.» Vedere una faccia che riconosceva sembrò soddisfare
Rogen, fece un cenno con la testa a Wayne ed entrambi misero la sicura
alle armi. «Cosa succede, sergente?»

«Ne sappiamo quanto voi, eravamo fuori per un'escursione nel luogo dov'è
precipitata quella nave, quando si è scatenato l'inferno. O i ruhar
hanno una visione molto elastica del loro accordo di cessate il fuoco, o
hanno deciso che vogliono riprendersi questo pianeta.» Ripensai a quello
che avevo sentito dire dal capitano Price. Non erano più informazioni
riservate: «La settimana scorsa, ho sentito che i kristang hanno
schierato la copertura aerea con l'eccezione di una singola fregata, in
cielo era in corso una specie di azione di flotta. Credo che i ruhar
abbiano attaccato Forte Freccia con un colpo di cannone a rotaia e
missili». La base del deposito di munizioni non si vedeva, una spalla
della catena montuosa chiudeva la visuale, come se non bastasse la nube
di fumo ancora denso. «Non abbiamo comunicazioni, impostate i vostri
zPhone in modalità di sola ricezione, così i ruhar non potranno usarli
per localizzarci.»

\hfill\break

«Devono sapere di questo posto», Rogen indicò le pesanti porte
all'ingresso del deposito di munizioni, «l'hanno costruito i criceti. E
vedranno che abbiamo aggiunto una baracca di guardia, così sapranno che
lo stiamo usando.» La baracca delle guardie era una struttura che, sulla
Terra, sarebbe sembrata un piccolo capanno degli attrezzi di quelli che
le persone hanno nei loro cortili, teneva le guardie al riparo dalle
piogge tropicali pomeridiane e basta.

«Non vedo un veicolo», notò Pope, «voi ragazzi avete un Crivee?»

«No», Wayne scosse la testa, «tra due ore i ragazzi che hanno il compito
di sostituirci dovrebbero arrivare con un camion che ci dovrebbe poi
riportare alla base.»

«Potete aprire quelle porte?» Indicai l'entrata del deposito di
munizioni.

«Abbiamo il codice», disse Rogan, «ma non siamo autorizzati a\ldots»

«Rogen, se hai tenuto d'occhio quelle armi per quando avrebbe piovuto
sul bagnato, ci siamo.» Sapeva cosa volevo dire, sebbene il cielo fosse
per lo più sgombro da nubi: «La base di Forte Freccia è stata colpita e,
per quanto ne sappiamo, siamo l'unica resistenza organizzata qui intorno
in questo momento. Abbiamo bisogno di armi. Apri quella porta».

\hfill\break

Eravamo fortunati, c'erano per lo più armi nel deposito di munizioni, ma
anche una piccola scorta di acqua e cibo che razziammo. Niente sale o
compresse di elettroliti, purtroppo, ma mi assicurai che tutti
mangiassero arachidi salate o brezel per reintegrare quello che avevamo
perso con il sudore. «Prendete tutti un M4 e munizioni, prendete
munizioni extra e quella roba buona con la punta esplosiva, non i
proiettili standard.» Tutti sapevano come identificare le munizioni
kristang. «Pope, Stallings, ehm, Newman, ehm, Wayne e tu», indicai i
cinque ragazzi più robusti del nostro gruppo, «prendete due Zinger
ciascuno. Tutti gli altri prendano uno Zinger e un At4, non sappiamo di
cosa avremo bisogno, quindi portiamo tutto.» Mi misi un paio di Zinger
sulle spalle per dare l'esempio. Erano pesanti, soprattutto sopra l'M4.
Pensandoci un attimo, mi tolsi la fondina dell'arma dal fianco e la
lasciai su uno scaffale. Non ha senso cercare di vuotare il mare con un
bicchiere.

Rogen aiutò Wayne a sollevare un Javelin: «Ehi, sei quel Joe Bishop
che\ldots».

«Sì, Rogen. Te ne parlerò dopo, promesso», ammesso che ci sarebbe stato,
un dopo.

«Giubbotto antiproiettile, sergente?», chiese Pope.

Mi accigliai: «Ehm, lascio fare a te», dissi, ma immediatamente capii
quanto fosse vigliacco da parte mia non prendere quella decisione.
«Aspetta, no. Niente giubbotto antiproiettile.» Questo contravveniva del
tutto alle regole dell'esercito: «Dobbiamo muoverci in fretta e siamo
già sovraccarichi». La gente annuì, mi sorprese, mi aspettavo resistenza
a quell'ordine. Tutti dovevano aver capito che il Kevlar non sarebbe
stato molto utile contro le armi di fanteria ruhar come i raggi di
particelle. Il giubbotto antiproiettile serviva soprattutto a proteggere
il busto di un soldato dai danni causati dalle schegge, non da colpi
diretti; il giubbotto antiproiettile mi aveva salvato da gravi ferite in
Nigeria, per quanto avessi risentito del suo peso nel caldo. Se fossi
sopravvissuto, mi sarei di sicuro preso una lavata di capo per aver
violato i regolamenti in combattimento. Se ci fosse stata ancora un'Unef
a redarguirmi.

Eravamo pronti, la gente mi guardava di nuovo. Cosa fare? Tornare
indietro di corsa lungo la strada di accesso alla base, in pieno sole,
era un'idea suicida. Se le cannoniere ruhar non si erano preoccupate di
noi prima, di sicuro lo avrebbero fatto quando ci saremmo avvicinati
alla base, portando delle armi. Il fatto di vedere aerei ruhar nel cielo
mi fece capire che non solo avevano colpito Forte Freccia, ma avevano
fatto sbarcare truppe per prendere il complesso della base del Vettore.
Questo significava che le truppe ruhar erano a terra e, se erano a
terra, avremmo potuto colpirle se ci fossimo avvicinati abbastanza.
«Qualcuno ha una mappa di questo posto, l'intero complesso del Vettore?»
Avevo giusto un'idea della zona nelle immediate vicinanze di Forte
Freccia, tutto qui. Da quando ero arrivato, il Vettore aveva inviato il
carico nello spazio solo tre volte, sempre quando ero fuori in servizio
di scorta. Il massimo che avevo visto e sentito del Vettore in azione
erano stati una scia di condensazione che sfrecciava nel cielo e un
basso rombo in lontananza.

\hfill\break

«Ce l'ho io, sergente», si offrì Amaro. Corse verso di me e mi porse il
suo zPhone. Lo presi e scorsi le piante che aveva scaricato,
rinfrescando la mia memoria annebbiata. Non fu di grande aiuto. Il tubo
di lancio aveva un unico tunnel di accesso parallelo per la manutenzione
sul lato sud e condotti laterali che si collegavano alla superficie
all'incirca ogni chilometro. Noi ci trovavamo sul lato sud, per arrivare
al condotto di accesso avremmo dovuto scalare la montagna, sopra il tubo
sepolto del Vettore. Non era un'opzione praticabile con cannoniere ruhar
che ronzavano in giro. Alzai gli occhi dallo schermo dello zPhone,
guardai i soldati che erano in attesa dei miei ordini. Avremmo potuto
cercare di tenere parte del Vettore, combattere una battaglia di
logoramento che i ruhar avrebbero vinto, mantenere quella posizione fino
a quando non fossimo tutti morti. Barattare vite con il tempo, in quella
che sapevo essere una battaglia senza speranza. Non sembrava un buon
piano. I ruhar avrebbero potuto limitarsi ad aspettarci fuori, o
impiegare un qualche tipo di gas e immobilizzarci. Esaminai di nuovo la
pianta, cercando di trovare ispirazione. In battaglia, le decisioni
dovevano essere prese su due piedi, e il mio cervello era già rallentato
dalla stanchezza. Vidi molti piccoli comparti fuori dal tunnel,
dov'erano alloggiate apparecchiature elettriche o d'altro tipo, erano
vicoli ciechi che sarebbero diventati trappole mortali se i ruhar
avessero sorpreso i nostri soldati lì dentro. «Aspetta, cos'è questo?»
Mostrai lo zPhone ad Amaro. C'era un altro condotto, un condotto
piccolo, parallelo a quello del Vettore, a nord rispetto a noi.

Lei strizzò gli occhi nella luce fioca, poi annunciò: «È il condotto per
il plasma dall'impianto di fusione, fornisce carburante per i magneti
del Vettore».

«Come lo sai?»

«Prima della guerra avevo in programma di diventare ingegnere
elettronico», disse. «M'interessava, feci un tour del Vettore la prima
volta che venni qui.»

A me nessuno aveva offerto un tour. «Questo condotto è pieno di plasma?
È un gas surriscaldato, giusto?»

«Qualcosa di simile, il plasma è un quarto stato della materia, non
gassoso, liquido o solido. Ma è maledettamente caldo.»

«Ancora?» Un'idea mi si stava formando nella testa. «Il Vettore non ha
sparato per quanto, tre giorni a oggi? Non immettono plasma nel condotto
a meno che non stiano caricando per un lancio, giusto?»

«Ah, sì, il prossimo lancio non è in programma per la prossima
settimana, perché non c'è una nave da carico ruhar che possa
raccoglierlo», mi guardò, con gli occhi spalancati. Doveva avere intuito
la mia folle idea. «È probabile che il condotto sia fresco ormai,
ma\ldots»

Mi guardai attorno nella vasta caverna che era stata scavata dai ruhar,
solo un angolino era utilizzato dall'Unef come deposito di munizioni.
L'estremità dello spazio era nascosta in un buio minaccioso, arrossato
dall'illuminazione di emergenza, perché l'elettricità era saltata subito
dopo che Forte Freccia era stato colpito. «Ci serviranno altre torce.»

Il condotto portava dritto nel cuore dell'impianto di fusione. Come il
tubo di lancio, aveva molti punti di accesso per la manutenzione. Il
portello più vicino era a meno di un chilometro di distanza risalendo la
strada, l'avevamo passato mentre andavamo al deposito di munizioni. Mi
misi alla guida della mia squadra improvvisata, tutti noi dovevamo
vedercela con il caldo, il carico extra di armi, la necessità di
sgattaiolare fuori dalla strada quando un aereo ruhar la sorvolava e lo
shock di una bella mattina trasformata in orrore. Quando raggiungemmo il
condotto di accesso, lasciai gli altri al sicuro fuori e portai Amaro
con me per esaminare il portello. C'erano ogni genere di controlli
elettrici e sensori, tutti disattivati in quel momento. E c'era una
ruota. Una semplice, grande ruota metallica. Appoggiai lo Javelin a
terra, afferrai la ruota e la feci girare. Una dozzina di giri e il
portello si stappò con uno sbuffo d'aria. Niente inferno di plasma a
bruciarmi i piedi. «Uhm», dissi infilando dentro la testa. Era un
condotto scuro che si estendeva per tutto il raggio del fascio di luce
della mia torcia. Circa tre metri di diametro. Spensi la torcia e
sbirciai nel buio. Nessuna luce splendeva dall'altra estremità. C'erano
pannelli luminosi sul soffitto del condotto, tutti spenti.

Amaro infilò la testa nel condotto: «Ehi. Eccolo qua. Più grande di
quanto pensassi. Anche qui è saltata la corrente, sergente, l'impianto
dev'essersi spento. O è stato colpito», ipotizzò Amaro.

«I ruhar devono essere stati attenti a non colpire il reattore, hanno
bisogno del Vettore intatto. È probabile che si sia spento in automatico
quando Forte Freccia è stato colpito.» La base era stata creata in una
parte della città costruita dai ruhar per i criceti che lavoravano al
complesso del Vettore, si trovava a nord del reattore a fusione. In
città vivevano ancora molti ruhar che facevano funzionare il Vettore e
il reattore, e si occupavano della loro manutenzione, i nostri tecnici
non erano in grado di fare nulla di utile con la complessa tecnologia
aliena. Una recinzione separava Forte Freccia dalla città, e i ruhar
erano confinati nelle loro aree. Sono sicuro che la flotta dei criceti
fosse stata attenta a non colpire la parte ruhar della città. Mi spinsi
fuori dal condotto e feci cenno agli altri di venire avanti. «Ecco il
piano, passeremo al contrattacco. Questo condotto porta all'impianto di
fusione, possiamo uscire lungo la strada. Lo seguiremo, sbucheremo nelle
retrovie del nemico e faremo un po' di casino. Andiamo a uccidere
qualche criceto questa mattina.»

\hfill\break

Per fortuna, nel condotto rimase buio pesto mentre lo attraversavamo.
Durante tutto il percorso ebbi la pelle d'oca, per paura che i ruhar ci
localizzassero e lanciassero qualche razzo o granata lungo il condotto,
o ci spazzassero via sfruttando lo spazio ristretto. O attivassero
qualunque meccanismo generasse il plasma, per immetterlo nel condotto,
cosa che ci avrebbe fritto tutti fino a farci diventare croccanti. In
segreto, speravo che il reattore fosse stato danneggiato, che non fosse
solo temporaneamente disattivato, perché il danneggiamento avrebbe fatto
sì che il plasma non potesse essere attivato da un momento all'altro.
Non c'erano telecamere in vista nel condotto perché la presenza di
plasma surriscaldato in funzione le avrebbe rese poco pratiche. Ci
dovevano essere dei sensori lungo la strada per monitorare il plasma, la
mia speranza era che non fossero in grado di rilevare la nostra
presenza. A gestire il Vettore e i macchinari del reattore erano i
ruhar, pensai che i sensori, compresi tutti i dispositivi di
sorveglianza, dovessero essere collegati a Forte Freccia da qualche
parte, sapevo che erano gli umani a gestire il sistema di sicurezza per
il complesso del Vettore e della base. Visto che Forte Freccia era stato
colpito, speravo che i ruhar non avessero accesso ai video di
sorveglianza perché, se così fosse stato, ci avrebbero individuati
appena saremmo usciti dal condotto ed entrati nell'area della base.

Quello con cui non avevo fatto i conti erano i numerosi componenti di
macchinari su cui saremmo inciampati. Nella mente, mi ero figurato il
condotto come un tubo liscio vuoto, un buco rotondo nel terreno. Non era
così. Amaro spiegò che il plasma era contenuto in un campo magnetico,
che richiedeva magneti disposti a un paio di metri di distanza l'uno
dall'altro. E i magneti richiedevano cavi di alimentazione che
necessitavano di una schermatura termica. Tutti quei componenti facevano
inciampare e cadere i soldati in corsa, che a loro volta facevano
inciampare e cadere quelli dietro. Dato che nel condotto si poteva stare
solo in fila indiana, chi cadeva non si procurava solo gomiti e
ginocchia ammaccati, ma faceva fermare l'intera colonna. La cosa che più
temevo era che dal fucile di qualcuno partisse un colpo accidentale, che
avrebbe messo in allarme i ruhar, perciò ordinai a tutti di avanzare a
passo svelto, ma accorto, e questo ridusse al minimo il rischio
d'inciampare. Fu un lungo viaggio, camminammo per chilometri sottoterra
nel buio claustrofobico, con una visuale immutabile, non fosse per i
numeri ruhar lungo il muro. Per fortuna, l'aria era fresca. Amaro disse
che era probabile si trattasse del freddo residuo, provocato dai magneti
superconduttori del tubo di lancio principale, che si diffondeva nel
condotto, rimasto senza plasma per giorni. A parte bestemmiare quando
cadevano o sbattevano contro i muri, i componenti della squadra
mantenevano la disciplina, parlando il meno possibile, in un sussurro.
Se fosse perché erano stati addestrati in modo ammirevole o perché erano
spaventati a morte come me, non avevo voglia di scoprirlo. Ero
terrorizzato al punto che la torcia mi tremava in mano e continuavo a
farla dondolare da un lato all'altro perché non si capisse.

Camminare per chilometri sottoterra nell'oscurità mi lasciava troppo
tempo per pensare. Stavo conducendo quelle persone a una morte inutile,
perché non riuscivo a pensare a qualcosa di meglio da fare? Qualcosa di
più intelligente? Il fatto era che non mi veniva in mente nient'altro da
fare. Eravamo soldati, eravamo stati attaccati, combattevamo. Semplice.
Se un ufficiale mi avesse ordinato di fare la stessa cosa che avevo
ordinato a quelle persone, avrei obbedito senza fare domande. Questo non
significava che fosse la cosa giusta o la cosa migliore da fare. Quello
che avevo detto a Collins era vero: se i ruhar si stavano riprendendo il
pianeta, presto avremmo potuto essere tutti morti. L'unica chance di
successo, di sopravvivenza, meglio, della missione era tenere impegnati
i ruhar con il combattimento a terra, con tattiche di guerriglia se
necessario, e guadagnare tempo per consentire ai kristang di riprendere
il controllo nello spazio. Se i kristang non fossero tornati o non
avessero potuto, saremmo stati spacciati comunque e avremmo dovuto
mettere a segno almeno un punto per l'umanità: farla pagare in qualche
modo ai criceti e mostrare ai kristang che gli umani potevano essere
utili, alleati affidabili, perché è di questo che la gente della Terra
aveva bisogno.

Quando ci avvicinammo finalmente all'estremità del condotto dove si
trovava la centrale elettrica, potei vedere davanti a noi dove il
condotto curvava e si ramificava verso il punto in cui, ipotizzai, era
generato il plasma. Ovunque fosse, non avevo intenzione di andarci. Mi
fermai e ordinai di spegnere le luci, mentre io e Amaro controllavamo la
pianta sul suo telefono. Cercai d'immaginare dove si trovasse il
reattore rispetto a Forte Freccia; a est, più in alto sulla montagna e
più vicino al tubo di lancio. Se eravamo vicini al punto in cui il
condotto del plasma si ramificava per connettersi al reattore, eravamo
andati troppo oltre. Chiedere a tutti di fare dietrofront negli angusti
confini del condotto e marciare a lungo a ritroso non era nelle mie
intenzioni. «Amaro, questo condotto finisce al reattore, o vicino?»

«No, il reattore fornisce energia al generatore di plasma. Il generatore
è la grande struttura rotonda sulla montagna, sembra una torre idrica o
una cisterna di petrolio. C'è un alto edificio bianco proprio accanto.
Dobbiamo essere nelle vicinanze di quello.»

«Ah, sì, va bene» Eravamo vicini a dove volevo essere, anche se avevo
dimenticato quanto in alto sulla montagna dovevamo trovarci allora.
«Pensavo fosse una cisterna d'acqua.» Mentre marciavamo, il condotto
saliva in modo così graduale che non me ne ero accorto. Il Vettore era
lungo molti chilometri e si arrampicava sulla montagna e la tagliava,
dalla base, a ovest di Forte Freccia fino all'estremità aperta del tubo
di lancio a est, su per l'altura e sull'altro lato della cresta. La
massa della montagna proteggeva Forte Freccia e il complesso del Vettore
dall'onda d'urto sonica che si generava quando le capsule di carico
lasciavano il tubo di lancio e colpivano l'atmosfera. In convoglio, se
stavamo per trovarci nell'impronta sonica di un lancio, dovevamo
fermarci un'ora prima, mettere in sicurezza i veicoli e indossare
protezioni per l'udito sotto i caschi. Fino ad allora, mi era successo
solo due volte durante il servizio di scorta: ci trovavamo a una
distanza tale da far sembrare il lancio una fioca striscia di luce nel
cielo, eppure avevo sentito il terreno rimbombare per le onde d'urto
ipersoniche. Tutti gli edifici all'interno e intorno al complesso del
Vettore, tra cui Forte Freccia, avevano sistemi d'insonorizzazione
installati nelle pareti e nelle finestre, e la base veniva isolata
durante ogni lancio. Un lancio era una cosa che avrei voluto vedere, ora
rischiavo di perdere quell'occasione.

Amaro e io ci spingemmo oltre la colonna e tornammo indietro di un
centinaio di metri, fino al punto in cui avevamo superato un portello
d'accesso. Anche i portelli avevano delle ruote all'interno, e quello in
particolare aveva una piccola e spessa finestra al centro. Avvicinare
l'occhio alla finestra era inutile, tutto quello che riuscivo a vedere
era il riflesso del mio bulbo oculare, dall'altra parte era buio pesto.
Se ci fossero stati dei criceti, saremmo stati fottuti. «Amaro, prepara
una granata.»

Amaro annuì cupa e io cominciai a girare piano piano la ruota. Non
cigolava, così la girai più veloce che potevo e aprii del tutto il
pesante portello.

Un corridoio buio. Un corridoio vuoto, lungo una trentina di metri, alla
cui estremità c'era una porta, una normale porta rettangolare
dall'aspetto solido. Feci cenno ad Amaro di seguirmi e agli altri di
restare fermi, e proseguii camminando piano. Sulla porta c'era un
tastierino spento, niente energia, e una leva al posto di una ruota. La
leva girò in modo sorprendentemente facile, la porta era
sorprendentemente pesante. Al di là si apriva l'interno di una sorta di
garage, con una saracinesca alta e larga e un camion ruhar parcheggiato
a sinistra. Era un semplice camion, non era stato convertito in Crivee.
È probabile che fosse stato utilizzato dai ruhar che si occupavano della
manutenzione delle macchine del Vettore. Il garage era caldo e umido,
c'era un condizionatore d'aria non funzionante incassato nel soffitto.
Amaro e io strisciammo attorno al camion, dove c'era una porta normale
che portava a una zona ufficio; due tavoli molto graffiati, un paio di
sedie usurate, un banco da lavoro, attrezzi, un bidone della spazzatura
con involucri di pranzi dei criceti. E polvere. Quel posto, supposi, non
era stato molto frequentato. Strisciammo con cautela nell'ufficio e
guardammo dalle finestre sudice.

C'era una splendida vista, in un certo senso. Eravamo in cima alla
montagna, con il complesso del Vettore, il reattore, la città, Forte
Freccia e il campo d'aviazione che si estendevano sotto di noi. C'era
un'altra porta che dava sull'esterno, di fronte alla grande saracinesca
si apriva uno spiazzo di cemento che connetteva alla strada d'accesso
sterrata più in là, si vedevano una tettoia e una gru sopra di esso. La
tettoia ci avrebbe fornito protezione dagli occhi che ci spiavano
dall'alto.

«Amaro, porta tutti qui. E chiudi il portello, in caso torni la
corrente. Un portello aperto potrebbe far scattare l'allarme.»

Uscii al riparo della tettoia. Sotto di me si estendeva in effetti Forte
Freccia. Quello che ne restava. Quando l'Unef lo aveva creato dalla
città dei criceti, avevamo abbattuto edifici per creare un perimetro e
installato una recinzione e un campo minato, perciò era facile
distinguere i confini della base. Al posto della sala mensa c'era un
cratere che emanava vapore, vapore invece che fumo, perché l'impattatore
del cannone a rotaia aveva danneggiato l'edificio ricreativo dall'altra
parte della strada e l'acqua della piscina incrinata si era riversata
nella voragine. I missili avevano anche colpito gli edifici delle
caserme, il che non aveva alcun senso poiché i ruhar dovevano sapere che
quelli sarebbero stati per lo più vuoti a metà pomeriggio, ma gli
edifici dell'amministrazione della base erano per lo più intatti, non
fosse per i danni causati da detriti volanti. Lo stesso valeva per la
maggior parte degli altri edifici all'interno della recinzione di Forte
Freccia.

Più in là, verso nord, il campo di aviazione era un disastro, sulle
piste si erano formati dei crateri, così non potevamo far decollare
nessun Dumbo, e ogni hangar aveva ricevuto un colpo diretto. Polli e
Poiane fracassati erano sparpagliati in giro per il campo, a occhio non
eravamo riusciti a far volare un solo uccello nell'aria, non fosse che
una colonna di fumo proveniente dalla giungla faceva pensare che potesse
salire da un aereo abbattuto. Era un disastro.

I miei uscivano in fila dalla porta per andare sotto la tettoia, io li
ammonii di restare nell'ombra e di non indossare occhiali da sole o
qualsiasi altra cosa potesse riflettere la luce stellare e tradire la
nostra posizione. La maggior parte di loro restò senza fiato quando vide
la distruzione, sentii molti dire in tono calmo che eravamo stati
fortunati a non trovarci in sala mensa.

La sala mensa. Quella che Forte Freccia aveva adibito a sala mensa era
stata usata con una funzione simile dai criceti prima di noi: stimai che
potesse agevolmente contenere quattrocento persone. Quattrocento, tutte
morte. Impossibile che qualcuno fosse uscito vivo da quel cratere che
emetteva vapore.

«Qual è il piano, sergente?», chiese Pope.

Il mio piano era quello di infiltrarci nella città attorno a Forte
Freccia e ripararci sotto gli edifici dei criceti per attaccare le forze
ruhar; immaginando che ci avrebbero combattuto casa per casa, piuttosto
che far saltare in aria metà della loro città per prenderci. Avremmo
potuto tenerli occupati per molte ore, anche giorni, dando ai kristang
il tempo di riprendere il controllo in cielo. Inoltre, finché gli umani
tenevano parte della città, i ruhar non avrebbero corso il rischio di
azionare il Vettore. Il motivo per cui avevo portato così tanti Zinger è
che non avremmo avuto alcuna possibilità di resistere se i ruhar
avessero potuto fluttuare nel cielo sopra di noi e colpirci con
precisione. Il semplice fatto di sparare un singolo Zinger ogni tanto
avrebbe costretto il supporto aereo ravvicinato ruhar a ritirarsi e
cambiare tattica. Lo sapevo, perché in Nigeria avevo visto che i
ribelli, sparando alla cieca contro i nostri elicotteri con lanciarazzi
non guidati, inducevano la nostra copertura aerea a precipitarsi al
riparo.

A nord, oltre il campo d'aviazione, due coppie di Avvoltoi stavano
volteggiando, una di esse forniva copertura, mentre l'altra sganciava
missili e mitragliava la giungla. Amaro indicò l'azione, alcuni umani
dovevano essere sopravvissuti e fuggiti dal campo d'aviazione per
combattere. Mentre guardavamo, un paio di Zinger schizzarono fuori dalla
giungla e sfrecciarono verso uno degli Avvoltoi; il primo missile fu
colpito da un fascio di particelle e andò fuori rotta; il secondo mancò
il bersaglio, ma gli esplose abbastanza vicino da far sì che l'Avvoltoio
barcollasse nell'aria e si dirigesse ondeggiando verso il campo
d'aviazione, lasciando una scia di fumo.

«Binocolo», ordinai, porgendo la mano a Pope. Feci un passo indietro
nell'ombra più profonda ed esaminai il campo d'aviazione. Diversi Dodo
erano in manutenzione sulla pista e un paio di Balene stavano
scaricando. I Dodo erano piccole navicelle da trasporto, mentre le
Balene erano enormi navicelle, più grandi persino dell'aereo Dumbo con
cui ero arrivato. Le Balene potevano trasportare grandi carichi giù
dall'orbita, avevano porte all'estremità posteriore e una rampa come un
aereo da carico. Una delle Balene all'aeroporto stava scaricando una
Poiana che aveva le ali piegate.

«Merda, non va bene», dissi in tono calmo, ma non abbastanza.

«Cosa?», chiese Pope.

«Se stanno già portando le Poiane, devono essere sicuri che resteranno
qui per un po'.» Come avremmo fatto ad arrivare in città senza essere
visti? Non avevo pensato che saremmo usciti dal condotto così in alto
sulla montagna, dovevamo essere a cinquecento metri sopra la città.
Perché non avevo capito che\ldots{}

«Arrivano!», gridò qualcuno dietro di me. Con un ruggito acuto delle
turbine, due navicelle raggiunsero la cresta dietro noi, volarono vicino
al garage o qualunque cosa fosse la struttura in cui avevamo trovato
riparo, e passarono di colpo in volo a punto fisso per atterrare sul
campo d'aviazione. Erano un Dodo scortato da un Avvoltoio. Ci
schiacciammo tutti contro la saracinesca, il più in fondo possibile
sotto la tettoia. Guardai mentre atterravano e mentre si assestavano,
una delle Balene decollò in una nuvola di polvere, prese quota e volò
vicino a noi, salendo sopra la cresta alle nostre spalle e guadagnando
velocità in un lampo. Vidi la sua pancia scorrere nel cielo mentre ci
sorvolava. La sua grande, grassa, pancia vulnerabile.

«Uhm.» Proprio così, ebbi un'intuizione: «Questa è la rotta stabilita
ora. Siamo proprio sotto la rotta stabilita». Ogni volta che ero stato a
Forte Freccia, avevo visto aerei avvicinarsi da est o da ovest, a volte
da nord, ma mai da sud, mai visti arrivare da oltre la montagna. Le cose
erano cambiate. «I nostri ragazzi laggiù, nella giungla, hanno fatto sì
che i ruhar temessero di arrivare da quella direzione. Merda! Stanno
volando proprio sopra di noi, non sanno che siamo qui!»

Con il binocolo, scrutai in direzione dell'altra Balena che stava ancora
scaricando. Alcuni carichi stavano scendendo dalla rampa posteriore e, a
giudicare dal numero di soldati che percorrevano le scale sul lato della
nave e si allineavano sulla pista, quella Balena aveva per lo più
trasportato truppe. Una Balena, secondo le informazioni fornite dai
kristang, poteva trasportare fino a seicento passeggeri. Vedendo
l'equipaggiamento che avevano addosso quelli sulla pista, pensai che un
aereo da trasporto truppe Balena non potesse contenere il numero massimo
di persone.

Gli Avvoltoi sopra la giungla oltre il campo d'aviazione sparavano
ancora un missile ogni tanto contro chiunque si trovasse laggiù. Era
molto probabile che, se fossimo riusciti ad abbattere anche solo una
navicella dalla nostra posizione, i ruhar avrebbero interrotto le
operazioni al campo d'aviazione, finché non fossero stati sicuri di aver
sradicato tutta la resistenza umana nella zona. Questo avrebbe ritardato
di molto il loro programma di occupazione del Vettore. Forse li avrebbe
pure costretti a deviare navi e truppe da altre aree e questo avrebbe
concesso una pausa alle nostre forze altrove su Paradiso. Anche se i
kristang non fossero tornati, o non fossero potuti tornare su Paradiso,
alla fine avrebbero di sicuro sentito parlare della posizione ferma
tenuta dall'Unef al Vettore, e rafforzare l'opinione dei kristang
sull'utilità degli umani in combattimento avrebbe aiutato gli abitanti
della Terra, che era lo scopo della Expeditionary Force delle Nazioni
unite. Le nostre forze non avevano alcuna possibilità di difendere la
Terra dai ruhar, avevamo bisogno che i kristang lo facessero per noi.
L'Unef era andata tra le stelle per fornire ai kristang un motivo per
preoccuparsi se la Terra fosse stata conquistata dai ruhar. In un certo
senso, la missione dell'Unef era SALVARE IL MONDO.

Tutte quelle maiuscole erano volute, per inciso, sembra molto più
drammatico così. Quando stai affrontando la forte possibilità di restare
intrappolato, per sempre, su un pianeta alieno controllato dal nemico,
avere un fondamento logico-drammatico per la tua missione ti fa sentire
meglio.

Ridiedi il binocolo a Pope: «Tu e, uhm, Wayne potete attraversare la
strada e andare sotto quegli alberi? Ho bisogno di un osservatore per
vedere cosa sta arrivando oltre il crinale dietro di noi. Con la mano
destra, alza le dita per dirci il numero; con la mano sinistra, ehm, un
pugno significa una Balena, il palmo aperto verso il basso è un Dodo e
il palmo aperto verso l'alto è un Avvoltoio, capito?». Non volevo
sprecare Zinger su un Avvoltoio, dovevamo avere un impatto più grande
che abbattere una cannoniera biposto.

Pope guardò a destra e a sinistra: «Sì, percorriamo la strada da questo
lato, sotto la copertura degli alberi, fino al punto in cui sporgono
oltre la strada, l'attraversiamo e torniamo indietro».

«Bene, lasciate Zinger e At4 qui. Appena spariamo gli Zinger, voi venite
qui di corsa e torniamo in quel condotto, per correre in fretta verso
l'uscita successiva.» I ruhar avrebbero fatto presto a trasformare il
garage in un mucchio di macerie dopo che avremmo sparato al loro aereo.
«Tutti gli altri, tenete pronti i vostri Zinger. Amaro, vedi se riesci
ad aprire questa saracinesca, dobbiamo uscire in fretta.» Mi tolsi di
dosso la tracolla di uno degli Zinger che avevo portato, aprii il
pannello di mira e premetti il primo pulsante per attivarlo. La guida a
infrarossi dell'arma sarebbe rimasta attiva per diverse ore.

Pope e Wayne appoggiarono per terra i loro missili, sbirciarono da sotto
il bordo della tettoia e sfrecciarono via sotto gli alberi. In un paio
di minuti, erano sotto un albero dall'altra parte della strada rispetto
a me, eravamo in grado di comunicare parlando ad alta voce, i segnali
con le mani sarebbero serviti quando gli aerei nemici avrebbero volato
con fragore sopra di noi. Non appena Pope e Wayne si furono sistemati
sotto gli alberi, sentimmo aerei in avvicinamento. Un Dodo e un
Avvoltoio.

Scossi la testa e feci un pollice verso con la mano. Se non fossero
comparsi presto bersagli più allettanti, mi sarei accontentato di un
Dodo; dovevo preoccuparmi che le truppe ruhar potessero percorrere la
strada di accesso e vederci. Di certo i ruhar avrebbero perlustrato
l'intera area.

Un altro aereo si avvicinò da dietro di noi. Un solo Avvoltoio. Di nuovo
feci il pollice verso. Controllando i miei, vidi che erano agitati:
dovevamo fare qualcosa, presto. «Controllate le sicure, gente, nessuno
spari finché non darò il segnale.» Li indussi uno a uno a comunicare
l'avvenuta ricezione del mio ordine, guardando ciascuno negli occhi. Per
quanto mi riguarda, cercai di mostrarmi calmo, perfino annoiato,
soffocai uno sbadiglio simulato. «Cazzo, vorrei che arrivassero presto,
ho fame», mi lamentai, e questo suscitò qualche risata nervosa.

La Balena nel campo di aviazione finì di scaricare i passeggeri e,
quando le truppe ruhar ebbero lasciato la pista a passo di marcia per
inoltrarsi nella giunga, decollò. Quella Balena volò quasi direttamente
sopra di noi, potevo leggere alcuni segni sul ventre e sulle ali. Mentre
si avvicinava, mi assicurai che tutti vedessero i miei pollici versi. Se
fossimo stati costretti, avrei attaccato una nave vuota sul volo di
ritorno, ma non era la mia preferenza.

Poi mi annoiai, un po' almeno. Dopo il decollo della Balena, passarono
forse dieci minuti senza alcuna attività di volo nella zona, a parte la
coppia di Avvoltoi sopra la giungla. Non stavano più mitragliando o
lanciando missili e non vedemmo neppure altri Zinger provenire dal
suolo. Durante la pausa nel traffico aereo, iniziai a chiedermi se
avessi fatto bene a non sparare a quel Dodo o alla Balena vuota che era
passata sopra di noi prima. Entrambe le azioni avrebbero raggiunto il
nostro obiettivo principale, quello d'interrompere le operazioni di volo
dei ruhar. In ogni caso, volevo colpire i criceti, colpirli duro,
fargliela pagare per tutte le persone che erano morte nella sala mensa
di Forte Freccia.

Mentre aspettavamo, mandai degli uomini a pattugliare la zona, dopo aver
innescato gli Zinger avremmo avuto bisogno di tornare subito di corsa
nel condotto e di non perdere tempo a raccogliere cose che avessimo
lasciato cadere. Merda, avrei dovuto lasciare il portello d'accesso al
condotto aperto, non perdere tempo a girare quella grossa ruota pesante.
Troppo tardi, avevo imparato la lezione.

«Balene!», gridò Wayne e agitò le braccia. Pope stava guardando nel
binocolo e agitò una mano eccitato. Wayne alzò un pugno e due dita, poi
palmo aperto a faccia in su con due dita. Due Balene, scortate da due
Avvoltoi. Non riuscivamo ancora a sentirli, dovevano essersi avvicinati
ad alta quota. Poi eccoli.

«Ci siamo, gente!», gridai. «Non sparate finché non do l'ordine! Abbiamo
due Balene; voi da questa parte prendete di mira la prima; voi da
quest'altra parte, voi siete con me sulla seconda. Restate sotto la
tettoia per ora.» Ero così agitato che mi costrinsi a controllare due
volte se il mio Zinger avesse la sicura. Mentre il suono urlante dei
motori a reazione delle Balene incrementava, mi avvicinai al bordo dello
spiazzo di cemento strisciando i piedi. Cosa sapevamo delle Balene?
Cercavo di ricordarlo. Erano ben protette per essere navicelle, con
molte torrette difensive a fascio di particelle. Con quindici di noi
pronti a sparare Zinger contro due soli bersagli, da distanza
ravvicinata, avevamo buone probabilità di fare centro. A bassa quota e a
bassa velocità speravo che una Balena colpita non avrebbe avuto il tempo
di riprendersi prima di sbattere contro il fianco della montagna, e
precipitare giù per il pendio.

Sporsi la testa fuori dalla protezione della tettoia per vedere dove
fossero le Balene. Erano vicine, sarebbero passate un po' a est rispetto
al punto in cui ci trovavamo, invece che sorvolarci. Nessuna di loro
aveva il sistema stealth in funzione. «Tenetevi pronti», ordinai a voce
alta e uscii dalla protezione della tettoia. «Guide a infrarossi!»

Con la sicura tolta, il display di mira acceso, centrai il reticolo di
puntamento sul retro della Balena in coda, sul motore a reazione
posteriore di babordo. Il motore era quasi del tutto inclinato verso il
basso, fornendo più portanza che spinta propulsiva, perché la Balena era
in fase di atterraggio. Il punto più vulnerabile del suo profilo di
volo. Lento e radente.

«Avete tutti agganciato un bersaglio?», gridai, poi ripetei la domanda.
Dovevamo sparare subito: con quindici di noi all'aperto, con gli Zinger
sulle spalle, eravamo certi che il paio di Avvoltoi di scorta ci avrebbe
presto avvistati. Quattordici persone confermarono di tenere saldamente
sotto tiro il bersaglio.

«Pronti, pronti, fuoco!» Sganciai il mio Zinger e altri quattordici
missili saltarono dai tubi di lancio. Lo Zinger veniva espulso dal
lanciatore monouso da un impulso magnetico, così che il getto di
propulsione non uccidesse l'utente, il razzo non si attivava finché non
si trovava a una distanza di cinquanta metri. Quando il propulsore del
razzo si attivò, persi quasi di vista i missili mentre andavano in
super-accelerazione. Tutto avvenne così in fretta che non so quanti dei
nostri Zinger siano stati deviati, neutralizzati o del tutto distrutti
dalle difese delle Balene. Quello che so è che le Balene barcollarono in
aria, mentre i motori a reazione anteriori e posteriori di babordo di
entrambe esplodevano e venivano divelti, e altri missili penetravano nel
ventre delle due sfortunate navi. Le Balene hanno motori a reazione nel
ventre e la capacità di restare in volo a punto fisso servendosi solo di
quelli, caratteristica che non fu loro d'aiuto quel giorno.

Ho avuto un incubo ricorrente di quel momento, ma penso, spero, sia un
falso ricordo; qualcosa che potrei non avere visto, qualcosa che
potrebbe non essere accaduto, ma che è incredibilmente nitido, marchiato
a fuoco nella mia memoria. Il compartimento passeggeri/merci di una
balena non ha finestre, le finestre sono punti deboli nella struttura di
un aereo, e la struttura di una Balena deve salire e scendere
dall'orbita centinaia, anche migliaia di volte nel corso della sua vita.
Ci sono piccole finestre nelle porte laterali, per qualche motivo ci
sono un paio di piccole finestre su entrambi i lati sul retro, vicino
alla rampa, e ci sono finestrini nella cabina di pilotaggio. I
finestrini della cabina di pilotaggio sono come quelli di un aereo
commerciale sulla Terra, abbastanza grandi perché i piloti vedano
all'esterno, non abbastanza perché qualcuno possa vedere da fuori gran
parte dell'interno. E le Balene dovevano trovarsi a ottocento metri,
forse, da noi, e un po' al di sopra della nostra posizione. Eppure, ho
un ricordo vivido e inquietante del pilota ruhar della seconda Balena,
quella cui avevo sparato io, che si gira per vedere da dove provenga la
minaccia missilistica e, proprio mentre il primo missile colpisce e fa
sbandare la prua della Balena nella mia direzione, per un breve secondo,
guarda dritto verso di me. Non solo verso la nostra posizione, non solo
verso un gruppo di soldati umani, non solo verso un particolare soldato
umano.

Verso di me.

Me, come se mi conoscesse, come se mi stesse chiedendo il perché, visto
che le nostre specie si erano sviluppate su diversi pianeti, a distanza
di migliaia di anni e avevamo vissuto vite del tutto diverse, perché le
nostre strade si erano incrociate così? Era necessario che l'unica volta
che ci eravamo incontrati in tutta la vita sparassi un missile alieno e
mettessi fine alla sua esistenza? Perché? Non stava interrogando
l'universo, non stava interrogando il karma o il destino o qualsiasi
essere divino adorasse, stava interrogando me.

Questo mi perseguitò per molto tempo.

Entrambe le Balene rotolarono sul fianco dopo avere perso i motori a
reazione, una di loro riuscì per un po' a conseguire una parvenza di
controllo, prima che entrambe precipitassero quasi in verticale nella
giungla sul pendio della montagna, ruzzolando ancora e ancora,
spaccandosi nel corso di esplosioni secondarie. Quello che avremmo
dovuto fare era ritirarci nel condotto subito dopo avere innescato i
missili, invece restammo a guardare, a bocca aperta o esultando, che è
quello che accade quando una squadra raffazzonata è guidata da un
sergente inesperto. Wayne mi riportò sul pezzo quando lui e Pope
attraversarono di corsa la strada per raggiungerci. «È ora di andare,
sergente?», chiese, a occhi spalancati.

«Sì. Sì, vai, Amaro, apri il portello!» Uno degli Avvoltoi si era
librato in verticale molto sopra di noi, l'altro aveva virato in cerchio
verso di noi e stava cercando un bersaglio da annientare. Stava cercando
noi. Senza pensarci, mi slacciai dalla spalla il secondo Zinger, felice
ora di avere portato quei trenta chili in più, agganciai l'Avvoltoio, e
premetti il grilletto. Le difese dell'Avvoltoio colpirono il mio Zinger
prima che si avvicinasse, ma raggiunsi comunque il mio obiettivo, perché
il mezzo virò e sfrecciò via ad alta velocità, concedendoci tempo
prezioso per rientrare nel condotto.

\hfill\break

Si scoprì che non ci sarebbe stato bisogno degli sforzi della nostra
squadra raffazzonata, che tutti i ruhar a bordo delle due Balene da noi
abbattute erano morti per niente. Neppure un'ora dopo, una task force
kristang tornò e le navi ruhar saltarono via, abbandonando le loro forze
su Paradiso. In quel momento, avevo ritirato la mia squadra nel
generatore di plasma, rifugiandomi in un posto che pensavo che i ruhar
non avrebbero rischiato di danneggiare. Dopo che le due Balene si erano
schiantate ed erano andate a fuoco e i ruhar sopravvissuti a terra
avevano superato lo shock iniziale, un paio di Avvoltoi iniziarono a
ronzare intorno come calabroni arrabbiati, coi pod delle armi in vista,
volando a bassa quota per mitragliare di tanto in tanto qualcosa a terra
coi laser. Sapevano dov'eravamo, potevo immaginare che i piloti degli
Avvoltoi avessero alzato la voce per ottenere l'autorizzazione a far
saltare in aria la mia squadra e avessero ricevuto ordine di non
rischiare di danneggiare l'indispensabile macchinario del Vettore.
Pensai che la nostra migliore chance di sopravvivenza fosse nasconderci
in un posto difendibile, che i ruhar non potessero bombardare con armi
pesanti, e guadagnare tempo. Il complesso del generatore di plasma
sembrava una scommessa vincente. Ci restavano solo munizioni da fucile,
At4, un paio di Zinger, qualche granata e la sensazione di avere fatto
del nostro meglio per vendicare i nostri morti. Trovammo una sala di
controllo ausiliaria che era sotterranea e non aveva finestre, perciò
dovemmo usare telecamere che i ruhar avevano stranamente dimenticato di
disattivare. Da una delle telecamere, vedemmo il cratere dove un tempo
si trovava la sala mensa, che ancora emanava vapore. Mi chiesi se le
persone che vi si trovavano avessero avuto sentore di guai prima che il
colpo di cannone a rotaia le vaporizzasse: avevano saputo, pochi secondi
prima dell'impatto, che le navi da guerra ruhar erano saltate in orbita?
Speravo proprio di no. Meglio, in quel caso, la beata ignoranza,
ascoltando un generale dell'Unef blaterare un discorso noioso,
fantasticando pensieri piacevoli, senza renderti conto di essere stato
ucciso finché non ti sei svegliato nell'Aldilà. Non avevano avuto la
stessa fortuna le persone che si trovavano al campo d'aviazione e al
parco macchine, che avevano avuto il tempo di vedere la sala mensa
scomparire e detriti piovere su di loro prima che missili intelligenti
sganciassero loro addosso grappoli di sub-munizioni e distruggessero
tutto ciò che poteva guidare o volare. Le cannoniere Avvoltoi che ci
avevano ignorato avevano mitragliato i sopravvissuti dispersi mentre
c'era molta confusione. L'Unef stimò seicentosettanta morti, dei circa
novecento umani di stanza a Forte Freccia quel giorno. Questa fu una
dura lezione per il futuro combattimento: se non controllavi la
posizione strategica eri bello che morto, e la posizione strategica in
questo caso era in orbita e sopra. Trincerarsi non aiutava molto quando
un colpo di cannone a rotaia poteva penetrare fino a trecento metri nel
terreno, se la nave che sparava incrementava seriamente la velocità
della volata. I ruhar avevano navi dedicate al bombardamento orbitale,
il pacchetto informativo che avevo visto raffigurava la lunga canna
sottile del cannone a rotaia, con reattori a fusione all'estremità e
nient'altro. Chi aveva bisogno di merdose testate nucleari quando un
dart di cannone a rotaia poteva fornire dieci chilotoni di offesa su un
bersaglio e le navi potevano scaricare un dart dopo l'altro, fino a
quando anche il bersaglio più difficile era una nube di atomi?

Considerando la natura del futuro combattimento, che diavolo ci facevano
gli umani lì?

Amaro fu la prima a notare il cambiamento della situazione, stava
monitorando una telecamera e mi chiamò quando vide che la coppia di
criceti che stava sorvegliando il nostro reattore a fusione aveva posato
le armi e si era alzata, risultando quindi in piena vista. Spostammo la
visualizzazione da una telecamera all'altra, la scena si ripeteva in
tutta la base: criceti che deponevano le armi e si allontanavano,
dirigendosi verso il campo d'aviazione. Avvoltoi atterravano, equipaggi
uscivano, lasciavano le porte aperte e se ne andavano. Avevamo avuto
giusto il tempo di fare ipotesi su quello che stava succedendo, quando
tutti i nostri zPhone squillarono all'unisono. La cavalleria era
tornata. I nostri alleati kristang avevano ripreso il comando dello
spazio attorno a Paradiso, e le forze ruhar a terra si erano arrese. Per
il momento.